\documentclass[11pt]{article}

% ---- theme variables (passed via pandoc -V) ----
\newcommand{\ThemeName}{DSA — Index}
\newcommand{\ThemeKicker}{DSA Track}
\newcommand{\ThemeNumber}{}

\usepackage{interviewstyle}
\setthemecolor{be123c}{fb7185}

% pandoc-required plumbing
\providecommand{\tightlist}{\setlength{\itemsep}{1pt}\setlength{\parskip}{0pt}}
\setlength{\emergencystretch}{3em}
\providecommand{\pandocbounded}[1]{#1}

% syntax-highlighting macros emitted by pandoc (skylighting)


\begin{document}

\makecoverpage

\thispagestyle{empty}
{\hypersetup{hidelinks}\tableofcontents}
\clearpage

The coding interview is the \#1 filter at FAANG. These files cover \textbf{pattern recognition}, \textbf{data structures}, and \textbf{algorithms} --- plus two \textbf{worked-solution} companions that solve every named interview question end-to-end (optimal algorithm, idea, runnable code, error-prone tips).

\begin{longtable}[]{@{}
  >{\raggedright\arraybackslash}p{(\columnwidth - 2\tabcolsep) * \real{0.1843}}
  >{\raggedright\arraybackslash}p{(\columnwidth - 2\tabcolsep) * \real{0.8157}}@{}}
\toprule\noalign{}
\rowcolor{acc!16}
\begin{minipage}[b]{\linewidth}\raggedright
File
\end{minipage} & \begin{minipage}[b]{\linewidth}\raggedright
What\textquotesingle s inside
\end{minipage} \\
\midrule\noalign{}
\endhead
\bottomrule\noalign{}
\endlastfoot
\url{01-patterns-and-templates.md} & \textbf{Start here.} 22 coding patterns with "when you see X \(\rightarrow\) use Y" triggers + copy-pasteable Python templates. The pattern-recognition cheat table is the core. \\
\rowcolor{zebra}
\url{02-data-structures.md} & Every DS from first principles: internals, ASCII diagrams, operation-complexity tables, when-to-use, language container mapping. \\
\url{03-algorithms-and-complexity.md} & Big-O \& Master Theorem, sorting/searching, all graph algorithms, DP families, greedy, backtracking, bit tricks, number theory. \\
\rowcolor{zebra}
\url{04-data-structure-problems.md} & \textbf{Worked solutions} to all 68 data-structure interview questions from file 02 --- arrays, strings, hash maps, linked lists, stacks, trees, heaps, tries, graphs, union-find, caches, segment trees. Each: optimal approach, idea, runnable Python, error-prone spots. \\
\url{05-algorithm-problems.md} & \textbf{Worked solutions} to all the algorithm questions from file 03 --- complexity Q\&A, sorting, binary search, graph algorithms, DP, backtracking, greedy. Same format. \\
\end{longtable}

\section{How to use}\label{how-to-use}

\begin{enumerate}
\def\labelenumi{\arabic{enumi}.}
\tightlist
\item
  \textbf{Learn the patterns} (file 01) --- recognition is 80\% of the battle.
\item
  \textbf{Drill the complexity tables} (files 02 + 03) --- interviewers always ask "what\textquotesingle s the time complexity?"
\item
  \textbf{Map every LeetCode problem you do back to a pattern} in file 01\textquotesingle s cheat table.
\end{enumerate}

\section{The universal coding-interview loop}\label{the-universal-coding-interview-loop}

\setcodelang{}

\begin{verbatim}
Clarify → Examples/edge cases → Brute force (state complexity)
   → Optimize (pick a pattern) → Code cleanly → Test/dry-run → State final complexity
\end{verbatim}

\begin{quote}
\textbf{Acceptable complexity vs N} (memorize): N\(\leq\)20 \(\rightarrow\) 2\^{}N/N! ok \(\cdot\) N\(\leq\)500 \(\rightarrow\) O(N\(^3\)) \(\cdot\) N\(\leq\)5000 \(\rightarrow\) O(N\(^2\)) \(\cdot\) N\(\leq\)10\(^6\) \(\rightarrow\) O(N log N) \(\cdot\) N\textgreater10\(^7\) \(\rightarrow\) O(N) or O(log N).
\end{quote}

\end{document}
